\section*{Leitidee}

Mit dem Begriff des \textbf{Vektors} wird Geometrie algebraisch zug�nglich.
Vektoren beschreiben nicht nur Punkte oder gerichtete Gr��en im Raum. Sie
k�nnen addiert und mit Zahlen multipliziert werden und bilden dadurch eine
\textbf{lineare Struktur}.

Aus diesen beiden grundlegenden Operationen entsteht das Prinzip der
\textbf{linearen Kombination}: Vektoren k�nnen miteinander kombiniert werden,
um neue Vektoren und damit neue Richtungen oder Punkte zu erzeugen. Dadurch
lassen sich Fragen danach untersuchen, welche Vektoren aus anderen aufgebaut
werden k�nnen, welche Richtungen voneinander abh�ngig sind und welche
Vektoren gemeinsam einen Raum aufspannen. Geometrische Fragestellungen werden
auf diese Weise mit algebraischen Methoden beschreibbar.

Diese Struktur ist nicht auf den zwei- oder dreidimensionalen reellen Raum
beschr�nkt. Vektoren k�nnen aus beliebig vielen reellen oder komplexen
Komponenten bestehen und damit Elemente von
\(\mathbb R^n\) beziehungsweise \(\mathbb C^n\) sein. Der Begriff des
\textbf{Vektorraums} fasst die daf�r wesentlichen algebraischen Eigenschaften
zusammen und l�st sie von einer bestimmten Dimension oder geometrischen
Darstellung.

Mit dem \textbf{Skalarprodukt} erh�lt die lineare Struktur zus�tzlich eine
geometrische Bedeutung. Es erm�glicht, L�ngen und Winkel zu beschreiben und
f�hrt insbesondere zum Prinzip der \textbf{orthogonalen Projektion}. Ein
Vektor kann dadurch auf eine Richtung oder einen Unterraum projiziert und
dort durch einen geeigneten Vektor bestm�glich approximiert werden.

Diese Idee reicht weit �ber geometrische Vektoren hinaus. Auch andere
mathematische Objekte, insbesondere Funktionen, k�nnen Vektorr�ume bilden.
Besitzen solche R�ume zus�tzlich ein geeignetes Skalarprodukt, lassen sich
Begriffe wie Orthogonalit�t, Projektion und Approximation auf Funktionen
�bertragen. Damit entsteht eine Verbindung zwischen der linearen Algebra und
sp�teren Methoden zur Darstellung und Approximation von Funktionen.

